\documentclass[11pt]{article}
\usepackage[margin=1in]{geometry}
\usepackage{amsmath,amssymb,amsthm,booktabs,longtable,hyperref}
\renewcommand{\labelitemi}{$\bullet$}
\newtheorem{theorem}{Theorem}
\newtheorem{proposition}[theorem]{Proposition}
\newtheorem{lemma}[theorem]{Lemma}
\newtheorem{corollary}[theorem]{Corollary}
\theoremstyle{remark}
\newtheorem{remark}[theorem]{Remark}
\newcommand{\Pois}{\mathrm{Pois}}
\newcommand{\Bin}{\mathrm{Bin}}
\title{Asymptotic i.i.d.\ prophet inequalities:\\ fixed-capacity limits and certified two-unit bounds}
\author{Anonymous}
\date{15 October 2026}
\begin{document}
\maketitle
\begin{center}
Version 0.1.0-candidate\\
Reserved DOI: \href{https://doi.org/10.5281/zenodo.23404235}{10.5281/zenodo.23404235}
\end{center}
\begin{abstract}
For every fixed positive integer $k$, we prove existence of the worst-case asymptotic ratio between the optimal online reward from at most $k$ selections among $n$ i.i.d.\ nonnegative observations and the expected sum of the $k$ largest observations. The observation law is known to the decision maker and may vary with $n$. The limit equals a continuum minimax value with prophet profile $P_k(x)=\mathbb E\min\{\Pois(x),k\}$. The proof combines compactness of cumulative acceptance functions, approximation at both quantile boundaries, and quantitative finite-horizon transfers in both directions. For $k=2$, explicit ball-arithmetic certificates give $0.83695\le\gamma_2^*\le0.8372914$ and strictly separate this limit from a previously proposed nonlinear-system value. The lower witness is a mixture of $36$ quantile policies; the upper witness is a measure with $41$ layer atoms, not an observation law. An occupancy-weighted construction, consistent with earlier finite-horizon flow formulations, shows that finite mixtures can be dominated by deterministic time- and state-dependent quantile thresholds. The numerical certificates are specific to $k=2$; no equality with the all-horizon infimum or attainment by a bounded threshold policy is claimed.
\end{abstract}

\section{Introduction}
Let $X_1,\ldots,X_n$ be independent with common known law $F$ on $[0,\infty)$ and $0<\mathbb E X_1<\infty$. A decision maker sees the observations sequentially and may irrevocably accept at most $k$ of them. For $n\ge k$, let $\mathrm{DP}_{n,k}(F)$ be the optimal expected accepted sum, and let $\mathrm{PR}_{n,k}(F)$ be the expected sum of the $k$ largest observations. We study
\[
 \gamma_{n,k}=\inf_F\frac{\mathrm{DP}_{n,k}(F)}{\mathrm{PR}_{n,k}(F)}.
\]
The infimum permits a different law at each horizon. Consequently, existence of a limit for each fixed law does not establish convergence of $\gamma_{n,k}$.

The main analytic result, Theorem~\ref{thm:general-k-game}, proves convergence for every fixed $k$ and identifies the limit with a continuum game. The compact strategy variable is the cumulative acceptance function, not a bounded set of thresholds or a compact family of observation laws. Lower transfer is uniform over laws for each fixed finite policy mixture. Upper transfer constructs a horizon-dependent finite-support law from each fixed step profile. Compactness and explicit origin and tail estimates connect these two directions. These arguments do not interchange a limit with an infimum over all laws.

\paragraph{Relation to earlier formulations.}
For $k=1$, the asymptotic constant and its optimal-policy interpretation are classical \cite{HK82,Kertz86,CFHOV21}. Jiang, Ma and Zhang~\cite[Definition~4.1, Eq.~(25), Theorem~4.2]{JMZ} give a finite-horizon semi-infinite linear programme for arbitrary $k,n$ using state occupancy and cumulative type coverage; their Proposition~4.8 controls quantile discretization at fixed $n$. Our cumulative acceptance formulation is closely related to that framework. Their Lemma~2.11 and Appendix~C also supply the relevant threshold-recovery precedent. We record an explicit continuous-time occupancy formula, not a new purification principle.

Brustle, Perez-Salazar and Verdugo~\cite[Theorems~1--2]{BPV} give an exact finite-horizon quantile formulation and an asymptotic lower bound
\[
 \gamma_{n,k}\ge\left(1-\frac{24k\log^2n}{n}\right)
                  \sum_{j=1}^k\theta_j^\star
\]
for sufficiently large $n$. This is a lower bound on a liminf, not an identification of the optimal limit. Molina, Gast, Loiseau and Perchet~\cite[Proposition~5.4, Eq.~(5.7)]{MGLP} analyze a related boundary-value construction; with equal selection and benchmark budgets it yields the same kind of lower bound. Their fixed-horizon finite-support reduction (Proposition~5.5) is distinct from convergence as the horizon grows. The asymptotic tightness question is also discussed in \cite{Brustle26}.

Quantile-kernel minimax and finite-horizon implementation are substantive methodological precedents in Loiseau et al.~\cite[Section~3, Appendix~C.4]{LMPPV}. Their principal model uses a single static threshold and a different endpoint topology. Here the controls depend on time and remaining budget, the benchmark is the top-$k$ sum, and compactness is imposed on cumulative coverage; both endpoint estimates are proved explicitly. Other quantile and restricted-policy formulations include \cite{PST,GMZ}, and related online-selection analyses include \cite{PSV}. Kong and Murthy~\cite{KM26} study fixed-law extreme-value asymptotics, which do not justify moving the infimum over laws through the limit. Cruz-Ossa, Perez-Salazar and Verdugo~\cite{CPV26} study single-threshold competition complexity with a potentially different observation budget. Neither comparison supplies the unrestricted adaptive limit proved here. A bounded source comparison found no earlier theorem with precisely this fixed-budget worst-case limit and characterization; this is not an exhaustive priority claim.

\paragraph{Certified two-unit consequence.}
The analytic theorem also makes finite continuum certificates sufficient for asymptotic bounds. The following numerical consequence uses the explicit witnesses and rigorous arithmetic described in Section~\ref{sec:certificates}.
\begin{theorem}[computer-assisted numerical bounds]\label{thm:main}
The limit $\gamma_2^*=\lim_{n\to\infty}\gamma_{n,2}$ exists, and
\[
 0.83695\le\gamma_2^*\le0.8372914<0.83730.
\]
Moreover
\[
 \theta_1^\star+\theta_2^\star
 <\frac{5723787}{6850000}<0.835590,\qquad
 \gamma_2^*-(\theta_1^\star+\theta_2^\star)
 >\frac{18641}{13700000}.
\]
\end{theorem}
Thus the nonlinear-system value is not the asymptotic optimum for two units. This does not settle a claim about $\inf_n\gamma_{n,2}$, or show where a particular restricted dual construction loses value. The strengthened certified bracket has width $0.0003414$; rounding the upper endpoint outward to $0.83730$ gives width $0.00035$. Numerical single-horizon bounds in \cite{JMZ,BPV} and fixed-law numerical asymptotics in \cite{KM26} are useful comparisons, but are not substitutes for this asymptotic certificate.

\section{A continuum characterization for every fixed number of selections}
\label{sec:fixed-k-extension}

Fix an integer $k\ge1$. For $n\ge k$, let $\gamma_{n,k}$ be the
infimum, over laws $F$ on $[0,\infty)$ with $0<\mathbb E_F X<\infty$,
of the optimal expected reward from at most $k$ online selections divided
by the expected sum of the $k$ largest of $n$ independent observations.
The restriction $n\ge k$ is immaterial for the limit; for smaller horizons
one may define the benchmark using the largest $\min(k,n)$ observations.
Set
\[
 P_k(x)=\mathbb E\min\{\mathrm{Pois}(x),k\}
       =\int_0^x\Pr(\mathrm{Pois}(s)\le k-1)\,ds.
\]
Then $P_k(0)=0$, $P_k'(0)=1$, $P_k(\infty)=k$, and
\[
 P_k''(x)=-e^{-x}\frac{x^{k-1}}{(k-1)!}\le0,
 \qquad 0<P_k(x)\le\min\{x,k\}\quad(x>0).
\]
In particular $P_k(x)/x\to1$ as $x\downarrow0$, including when $k=1$.

A nonnegative bounded measurable threshold vector $a=(a_1,\ldots,a_k)$ accepts an
arrival of scaled tail quantile $s$ at time $t$ with $j$ slots remaining
when $s\le a_j(t)$. Use column state probabilities indexed by
$0,\ldots,k$, initially $\pi(0)=e_k$, with generator
\[
 Q_{j-1,j}(t)=a_j(t),\qquad Q_{j,j}(t)=-a_j(t)\quad(j\ge1),
 \qquad Q_{\cdot,0}=0.
\]
The state equation is $\pi'=Q\pi$. Define
\[
 W^a(x)=\int_0^1\sum_{j=1}^k\pi_j^a(t)\min\{x,a_j(t)\}\,dt.
\]
An equal-mesh threshold is constant on each interval of some finite
equal partition of $[0,1]$. Let $\mathcal C_k$ be the pointwise closure
on $(0,\infty)$ of the convex hull of their functions $W^a$.

A profile $v$ is a nonzero nonnegative nonincreasing step function of
finite height, supported in $(0,S]$ for some finite $S$. Write
$M_v=\int_0^S v(s)\,ds$, let $v_1$ be its maximum height, and put
\[
 \Psi_v(z)=\int_0^S(v(s)-z)^+\,ds\quad(z\in\mathbb R).
\]
Let $V_0^v=0$ and define $V_j^v$, $1\le j\le k$, by
\begin{equation}\label{eq:general-k-values}
 (V_j^v)'(t)=-\Psi_v(V_j^v(t)-V_{j-1}^v(t)),
 \qquad V_j^v(1)=0.
\end{equation}

\begin{theorem}\label{thm:general-k-game}
For every fixed integer $k\ge1$, the limit of $\gamma_{n,k}$ exists and
\[
 \lim_{n\to\infty}\gamma_{n,k}
 =\inf_{v}\frac{V_k^v(0)}{\int v\,dP_k}
 =\max_{W\in\mathcal C_k}\inf_{x>0}\frac{W(x)}{P_k(x)}.
\]
The infimum is over the profiles just defined. The maximum is attained
in $\mathcal C_k$; no attainment by a bounded threshold policy or by a
profile is asserted. This statement does not identify the limit with
$\inf_n\gamma_{n,k}$ or give a rate uniform in $k$.
\end{theorem}

\subsection{Finite-horizon lower transfer}

For $n\ge k$ and $0\le x\le n$, define
\[
 P_{n,k}(x)=\int_0^x\Pr(\mathrm{Bin}(n-1,s/n)\le k-1)\,ds.
\]
Summing the beta integrals for ranks $0,\ldots,k-1$ gives
\begin{equation}\label{eq:general-k-finite-tail}
 P_{n,k}(n)=n\sum_{r=0}^{k-1}\binom{n-1}{r}
 \int_0^1q^r(1-q)^{n-1-r}\,dq=k.
\end{equation}
Let $a$ have $T$ time pieces and threshold bound $A$. For
$n\ge\max\{k,2A\}$, at step $i$ use $t_i=(i-1)/n$ and accept iff
$nU_i\le a_j(t_i)$ with $j\ge1$ slots remaining. With $p_{j,i}$ the
pre-step state probabilities, set
\[
 W_n^a(x)=\frac1n\sum_{i=1}^n\sum_{j=1}^k
 p_{j,i}\min\{x,a_j(t_i)\}.
\]
For a finite mixture write $\overline W_n$ and $\overline W$ for the
weighted cumulative payoffs, and put
\[
g_n=\inf_{0<x\le n}\overline W_n(x)/P_{n,k}(x).
\]

\begin{lemma}\label{lem:general-k-lower-transfer}
Every such mixture guarantees
$\mathrm{DP}_{n,k}(F)\ge g_n\mathrm{PR}_{n,k}(F)$ for every admissible
$F$, and
\[
 \liminf_{n\to\infty}g_n\ge
 c:=\inf_{x>0}\frac{\overline W(x)}{P_k(x)}.
\]
For a single policy, writing
\[
 D=e(4A^2+4AT)+2A+2T,\qquad \delta_n=D/n,
\]
one has $|W_n^a(x)-W^a(x)|\le\delta_n x$ and
\begin{equation}\label{eq:general-k-benchmark-error}
 |P_{n,k}(x)-P_k(x)|\le\frac{x^3/3+x^2/2}{n}.
\end{equation}
For a fixed finite mixture use the maxima of its component bounds.
\end{lemma}

\begin{proof}
Let $\psi(q)=\inf\{z:F(z)\ge1-q\}$ and realize
$X_i=\psi(U_i)$ with independent uniform $U_i$. Observed atoms are
handled by drawing the conditional quantile using fresh randomness;
this gives admissible randomized decisions and the stated iid joint
realization. Put $v_n(s)=\psi(s/n)$. The pre-step state is independent
of $U_i$, so the policy reward is $\int_0^n v_n\,dW_n^a$. An item at
quantile $q$ is among the top $k$ precisely when at most $k-1$ other
uniform quantiles are smaller. Thus
$\mathrm{PR}_{n,k}(F)=\int_0^n v_n\,dP_{n,k}$. Ties among observed
values are harmless. Every strict level set of nonincreasing $v_n$ is
an initial interval; Tonelli's layer-cake identity therefore turns
$\overline W_n\ge g_nP_{n,k}$ into the claimed reward inequality.
Finite mean makes both rewards finite. Mixtures are selected before
arrivals, and the same argument applies to their weighted payoffs.

In the induced $\ell_1$ matrix norm, $\|Q(t)\|\le2A$, independent of
the number of states. The discrete kernel is $K_i=I+Q(t_i)/n$. If
there is no threshold breakpoint inside its cell, its exact continuous
propagator $E_i$ satisfies
\[
 \|E_i-K_i\|\le(2A/n)^2e^{2A/n}.
\]
At most $T-1$ cells contain a breakpoint; on these the bound
$4Ae^{2A/n}/n$ suffices. All propagators are stochastic and hence
contractive in $\ell_1$. Telescoping gives
\[
 \max_i\|p_{\cdot,i}-\pi(t_i)\|_1
 \le e(4A^2+4AT)/n.
\]
Within a cell $\|\pi(t)-\pi(t_i)\|_1\le2A(t-t_i)$.
The integrand entries $\min\{x,a_j\}$ are bounded by $x$;
separating the breakpoint cells in the Riemann sum gives the displayed
$\delta_nx$ bound.

For completeness,
$d_{\rm TV}(\mathrm{Bern}(p),\mathrm{Pois}(p))
=p(1-e^{-p})\le p^2$ for $0\le p\le1$.
Coupling independent summands and then Poisson variables differing
in mean by $s/n$ gives
\[
 d_{\rm TV}(\mathrm{Bin}(n-1,s/n),\mathrm{Pois}(s))
 \le(s^2+s)/n.
\]
Apply this to the event $\{0,\ldots,k-1\}$ and integrate to obtain
\eqref{eq:general-k-benchmark-error}. Mixture errors are convex
combinations of the component errors.

If $c=0$, nonnegativity of $g_n$ proves the conclusion directly,
including the case $A=0$. Suppose $c>0$, so $A>0$, and set
$\kappa=P_k(A)/A>0$. For sufficiently large $n$,
$c-\delta_n/\kappa>0$. Concavity gives $P_k(x)\ge\kappa x$ on
$(0,A]$, and hence
\[
 \frac{\overline W_n(x)}{P_{n,k}(x)}
 \ge\frac{c-\delta_n/\kappa}
 {1+(A^2/3+A/2)/(\kappa n)}.
\]
The positive numerator justifies replacing the denominator by its
upper bound. For $A\le x\le n$, the numerator is constant and
$P_{n,k}(x)\le k$. Since
$\overline W(A)=\overline W(\infty)\ge kc$,
\[
 \frac{\overline W_n(x)}{P_{n,k}(x)}\ge c-\delta_nA/k.
\]
Both lower bounds tend to $c$.
\end{proof}

\subsection{Verification and quantitative upper transfer}

\begin{lemma}\label{lem:general-k-verification}
The system \eqref{eq:general-k-values} has a unique solution and
$V_j\ge V_{j-1}\ge0$ for all $1\le j\le k$. For uniformly bounded
Borel acceptance sets $\mathcal A_j(t)$ with jointly Borel sets
$\{(t,s):s\in\mathcal A_j(t)\}$, the associated continuum
policy reward is at most $V_k(0)$. Equality holds for
\[
 \mathcal A_j^*(t)=\{s:v(s)>V_j(t)-V_{j-1}(t)\}.
\]
If $\Phi_0=0$ and continuous piecewise absolutely continuous functions
$\Phi_j$ satisfy $\Phi_j(1)\ge0$ and
$\Phi_j'\le-\Psi_v(\Phi_j-\Phi_{j-1})$ almost everywhere, then
$V_j\le\Phi_j$ for every $j$.
\end{lemma}

\begin{proof}
The function $\Psi_v$ is nonnegative, nonincreasing and globally
$S$-Lipschitz on $\mathbb R$, so the finite-dimensional vector field
is globally Lipschitz and the terminal solution exists uniquely.
First $V_1\ge0$. Inductively let $D_j=V_j-V_{j-1}$ and suppose
$D_{j-1}\ge0$. On an interval where $D_j<0$,
\[
 D_j'=\Psi_v(D_{j-1})-\Psi_v(D_j)\le0.
\]
Such a negative interval cannot return to the terminal value
$D_j(1)=0$, proving $D_j\ge0$.

Write $b_j=|\mathcal A_j|$,
$r_j=\int_{\mathcal A_j}v$, and $\Delta_j=V_j-V_{j-1}\ge0$.
Then $\Psi_v(\Delta_j)\ge r_j-b_j\Delta_j$, even if acceptance
extends beyond the support of $v$. For the corresponding state law,
\[
 \frac{d}{dt}\sum_{j=0}^k\pi_jV_j
 =\sum_{j=1}^k\pi_j[-\Psi_v(\Delta_j)-b_j\Delta_j]
 \le-\sum_{j=1}^k\pi_jr_j.
\]
Integration from 0 to 1 gives $J\le V_k(0)$. For the displayed
maximizing sets, $\Psi_v(\Delta_j)=r_j-b_j\Delta_j$, giving equality.
These sets are bounded by $S$, Borel, and initial intervals up to
irrelevant endpoints, hence are bounded measurable quantile thresholds.

For the final assertion, proceed by induction in $j$. If
$V_{j-1}\le\Phi_{j-1}$ then
$V_j'\ge-\Psi_v(V_j-\Phi_{j-1})$. On a positive component of
$V_j-\Phi_j$, its derivative is nonnegative by monotonicity of
$\Psi_v$, so it cannot return to its nonpositive terminal value.
Thus that positive component cannot exist.
\end{proof}

\begin{proposition}\label{prop:general-k-upper-transfer}
Let $v=v_m$ on $(e_{m-1},e_m]$, with $e_M=S$ and $v=0$ beyond $S$.
For $n\ge\max\{k,2S\}$, let $F_n$ put mass
$(e_m-e_{m-1})/n$ at $v_m$ and the remaining mass at zero. Then
\[
 \mathrm{DP}_{n,k}(F_n)
 \le V_k(0)+\frac{(2eS^2+S)M_v}{n},
\]
\[
 \mathrm{PR}_{n,k}(F_n)
 \ge\int v\,dP_k-\frac{v_1(S^3/3+S^2/2)}n.
\]
Consequently $\limsup_n\gamma_{n,k}\le V_k(0)/\int v\,dP_k$.
\end{proposition}

\begin{proof}
Backward induction gives a deterministic optimal finite-horizon policy
whose decisions depend on time, remaining slots and current value.
It may reject zero without loss. In quantile coordinates its acceptance
sets $\mathcal A_{i,j}$ lie in $(0,S]$. Put
$\alpha_{i,j}=\int_{\mathcal A_{i,j}}v\le M_v$; its reward is
$n^{-1}\sum_{i,j}p_{j,i}\alpha_{i,j}$.

Use the same sets constantly on the corresponding continuum cells.
Their generators have norm at most $2S$, so
\[
 \|E_i-K_i\|\le2eS^2/n^2,
 \quad \max_i\|p_{\cdot,i}-\pi(t_i)\|_1\le2eS^2/n,
 \quad \|\pi(t)-\pi(t_i)\|_1\le2S(t-t_i).
\]
Multiplying by the reward-vector bound $M_v$ and summing cells gives
\[
 n^{-1}\sum_{i,j}p_{j,i}\alpha_{i,j}
 \le\int_0^1\sum_j\pi_j(t)\alpha_j(t)\,dt
       +(2eS^2+S)M_v/n.
\]
The integral is at most $V_k(0)$ by
Lemma~\ref{lem:general-k-verification}. The prophet identity and
layer cake, using \eqref{eq:general-k-benchmark-error} on $[0,S]$,
give the denominator bound. Its limit is positive because $v$ is
nonzero. Divide for sufficiently large $n$ and use
$\gamma_{n,k}\le\mathrm{DP}_{n,k}(F_n)/\mathrm{PR}_{n,k}(F_n)$.
\end{proof}

For numerical certification at any fixed $k$, the verification lemma
permits a conservative backward grid construction, computed in
increasing state order at each time step:
\[
 \Phi_j(t_i)\ge\Phi_j(t_{i+1})
 +h\Psi_v\bigl(\Phi_j(t_{i+1})-\Phi_{j-1}(t_i)\bigr),
 \qquad \Phi_j(1)=0,
\]
with upward rounding. Increments are nonnegative. Linear interpolation
satisfies
$\Phi_j(t)-\Phi_{j-1}(t)\ge
\Phi_j(t_{i+1})-\Phi_{j-1}(t_i)$, which proves the required slope
inequality. This is a valid general scheme, not a new numerical witness.

\subsection{Compact minimax and control of both boundaries}

\begin{proof}[Proof of Theorem~\ref{thm:general-k-game}]
Let $c$ denote the infimum over profiles in the theorem. The upper
transfer gives $\limsup_n\gamma_{n,k}\le c$. Clearly $c\ge0$;
also $c\le1$, since for $v=\mathbf1_{(0,x]}$ both the optimal reward
and prophet reward are $P_k(x)$, obtained by taking up to $k$ positive
Poisson arrivals.

Let $\mathcal K_k$ be the set of nondecreasing concave functions
$W:(0,\infty)\to[0,k]$ with $W(x)\le x$. Every $W^a$ lies in it:
its slope is at most $\int_0^1\sum_{j\ge1}\pi_j\le1$, and its
limit at infinity is the expected number of acceptances, at most $k$.
The defining inequalities of $\mathcal K_k$ are closed finite-evaluation
conditions. Thus it is a convex closed subset of the compact product
$[0,k]^{(0,\infty)}$, and $\mathcal C_k$ is compact and convex.
Every member extends by $W(0)=0$ and is $1$-Lipschitz: for $0<x<y$,
concavity gives
\[
 0\le\frac{W(y)-W(x)}{y-x}\le\frac{W(x)}x\le1.
\]

For each profile, the optimal bounded measurable thresholds from
Lemma~\ref{lem:general-k-verification} can be approximated in $L^1$
by equal-mesh thresholds with values in $[0,S]$. Indeed the cell
averages converge in $L^1$. For
$\|a-b\|_{L^1}=\int_0^1\sum_{j=1}^k|a_j-b_j|$, stochastic
contraction and Duhamel's formula give
\[
 \sup_t\|\pi^a(t)-\pi^b(t)\|_1\le2\|a-b\|_{L^1}.
\]
The reward integrals differ by at most $v_1|a_j-b_j|$ in state $j$,
so
\[
 |J_v(a)-J_v(b)|\le(2M_v+v_1)\|a-b\|_{L^1}.
\]
Consequently the supremum of the profile reward over equal-mesh
policies is $V_k^v(0)$.

Write profiles as $v_\nu(s)=\nu((s,\infty))$, where $\nu$ is a
nonzero finitely supported positive measure on $(0,\infty)$.
Fubini gives $J_{v_\nu}(a)=\int W^a\,d\nu$ and
$\int v_\nu\,dP_k=\int P_k\,d\nu$. Values at profile endpoints
affect none of these expressions or the ODE. Homogeneity of
\eqref{eq:general-k-values} permits normalization
$\int P_k\,d\nu=1$; let $\mathcal N_k$ be this convex set of
normalized finite measures. Finite evaluation makes
$W\mapsto\int W\,d\nu$ continuous and affine. It has the same
supremum over the generating policies as over their convex hull and
closure, so compactness gives
\[
 \max_{W\in\mathcal C_k}\int W\,d\nu=V_k^{v_\nu}(0),
 \qquad c=\inf_{\nu\in\mathcal N_k}V_k^{v_\nu}(0).
\]

We give the compact minimax argument explicitly. For each $x>0$ set
$H_x=\{W\in\mathcal C_k:W(x)\ge cP_k(x)\}$. If a finite
subfamily indexed by $x_1,\ldots,x_m$ had empty intersection, the
compact convex set of vectors
\[
 \left\{\left(\frac{W(x_\ell)}{P_k(x_\ell)}\right)_{\ell=1}^m:
                    W\in\mathcal C_k\right\}
\]
would be disjoint from $[c,\infty)^m$. Strict finite-dimensional
separation supplies nonnegative weights $\lambda_\ell$ with
$\sum_\ell\lambda_\ell=1$ such that
\[
 \max_{W\in\mathcal C_k}\sum_\ell
 \lambda_\ell\frac{W(x_\ell)}{P_k(x_\ell)}<c.
\]
The separating weights are nonnegative because the upper orthant is
unbounded in each positive coordinate direction. But
$\nu=\sum_\ell\lambda_\ell\delta_{x_\ell}/P_k(x_\ell)$ lies
in $\mathcal N_k$, contradicting the definition of $c$ and the
preceding payoff equality. The closed sets $H_x$ therefore have the
finite-intersection property, and compactness supplies
$W^*\in\mathcal C_k$ with $W^*\ge cP_k$ everywhere. Conversely,
for every $W\in\mathcal C_k$ and $\nu\in\mathcal N_k$,
\[
 \inf_{x>0}\frac{W(x)}{P_k(x)}
 \le\int W\,d\nu\le V_k^{v_\nu}(0).
\]
Taking the infimum over $\nu$ proves the minimax equality and
attainment in $\mathcal C_k$ without requiring $\mathcal N_k$ compact.

If $c=0$, nonnegativity and the upper transfer already prove limit
existence. Suppose $c>0$ and choose
$0<\varepsilon<\min\{c,1/2\}$. Choose $0<\delta<1<L$ with
\[
 P_k(\delta)/\delta\ge1-\varepsilon,
 \qquad P_k(L)\ge k(1-\varepsilon).
\]
Pointwise closure and the common $1$-Lipschitz bound give a finite
convex combination $W$ of equal-mesh policy payoffs with
\[
 \sup_{x\in[\delta,L]}|W(x)-W^*(x)|\le\varepsilon P_k(\delta).
\]
To see this directly, approximate at finitely many sufficiently fine
grid points, then apply the Lipschitz bounds between them. Thus
$W/P_k\ge c-\varepsilon$ on $[\delta,L]$. At the lower boundary,
concavity and $P_k(x)\le x$ give, for $0<x<\delta$,
\[
 \frac{W(x)}{P_k(x)}\ge\frac{W(\delta)}\delta
 \ge(c-\varepsilon)(1-\varepsilon).
\]
At the upper boundary, monotonicity and $P_k(x)\le k$ give, for $x>L$,
\[
 \frac{W(x)}{P_k(x)}\ge\frac{W(L)}k
 \ge(c-\varepsilon)\frac{P_k(L)}k
 \ge(c-\varepsilon)(1-\varepsilon).
\]
The fixed finite mixture representing $W$ is eligible for
Lemma~\ref{lem:general-k-lower-transfer}; hence
\[
 \liminf_n\gamma_{n,k}\ge(c-\varepsilon)(1-\varepsilon).
\]
Let $\varepsilon\downarrow0$ and combine with the upper transfer.
The order of choices is $k$, then $\varepsilon$, then the finite
mixture, then sufficiently large $n$. No joint uniform convergence
over all distributions or policies has been assumed.
\end{proof}

\subsection{Occupancy-weighted domination of finite mixtures}

The following constructive consequence is naturally expressed in the
state occupancy and acceptance flow framework of Jiang, Ma and
Zhang~\cite{JMZ}. We do not claim a novel purification principle.
It clarifies the role of mixtures as computational representations.

\begin{proposition}\label{prop:general-k-occupancy-domination}
Every finite mixture of bounded measurable quantile policies is
dominated, simultaneously for all $x>0$, by one bounded measurable
time- and state-dependent quantile policy. The corresponding statement
holds at each finite horizon for feasible quantile thresholds.
\end{proposition}

\begin{proof}
Let the mixture weights be $\mu_r$, thresholds $a_{r,j}$ and state
probabilities $\pi_{r,j}$. Put
\[
 p_j(t)=\sum_r\mu_r\pi_{r,j}(t),\qquad
 b_j(t)=
 \begin{cases}
 \displaystyle\frac{\sum_r\mu_r\pi_{r,j}(t)a_{r,j}(t)}{p_j(t)},
       &p_j(t)>0,\\
 0,&p_j(t)=0.
 \end{cases}
\]
These are measurable thresholds bounded by the maximum component
threshold. If $p_j=0$ all the nonnegative weighted occupancy terms
vanish, so $p_jb_j=\sum_r\mu_r\pi_{r,j}a_{r,j}$ in all cases.
Writing $p_{k+1}b_{k+1}=0$ and $b_0=0$, the averaged state equations
are
\[
 p_j'=p_{j+1}b_{j+1}-p_jb_j\quad(0\le j\le k),
 \qquad p(0)=e_k.
\]
Uniqueness for this bounded measurable linear system identifies $p$
with the state probabilities of policy $b$. Concavity of
$u\mapsto\min\{x,u\}$ gives
\[
 p_j(t)\min\{x,b_j(t)\}
 \ge\sum_r\mu_r\pi_{r,j}(t)\min\{x,a_{r,j}(t)\}.
\]
Summing and integrating proves $W^b\ge\sum_r\mu_rW^{a_r}$.
At a finite horizon define $b_{j,i}$ by the same weighted formula
using pre-step state probabilities. The acceptance flows agree, so
the state probabilities agree by induction. The same concavity
inequality proves domination of $W_n$ and hence, by layer cake,
reward domination for every admissible law.
\end{proof}

The dominating continuum thresholds need not retain an equal mesh.
The proposition implies equality of the game value with the
\emph{supremum} guarantee over bounded measurable deterministic
quantile thresholds, because each approximating finite mixture can
be dominated. Conversely, every bounded measurable threshold payoff
belongs to $\mathcal C_k$: approximate its thresholds in $L^1$ by
equal-mesh cell averages and apply the reward-continuity estimate above
to each indicator profile $v=\mathbf1_{(0,x]}$. This gives pointwise
convergence of the cumulative payoffs. The maximum over $\mathcal C_k$
therefore also bounds the deterministic supremum from above.
This does not prove attainment of that supremum, nor
remove the separate quantile randomization needed for atoms in $F$.
The analytic results of this section provide no numerical bracket or
nonlinear-system separation for $k\ge3$; such conclusions require
their own witnesses and certificates.

\section{Certified two-unit bounds}\label{sec:certificates}
In this section $k=2$, $P=P_2=2-e^{-x}(2+x)$, and $V_j=V_j^v$.
Thus $P''(x)=-xe^{-x}\le0$. Write $\bar W$ for the weighted cumulative
acceptance function of a finite mixture. The next two propositions
certify particular rational inputs; they do not certify the search
procedure or its optimality.

\begin{proposition}[computer-assisted]\label{prop:lowcert}
For the explicit mixture of $36$ rational quantile policies with $T=64$ in the data file \texttt{mix\_fine.json} (rational weights summing to $1$), $\bar W(x)\ge 0.83695\,P(x)$ for all $x\ge0$.
\end{proposition}
\begin{proof}[Proof (certificate)]
On each time piece the chain of each policy has a constant generator $M$; both $\pi$ at the end of the piece and $\int\pi\,dt$ over the piece are read off from $\exp\bigl(h\begin{smallmatrix}M & I\\ 0&0\end{smallmatrix}\bigr)$, $h=1/64$, computed as a rigorous ball. This gives $\bar W(x)=\sum_bw_b\min\{x,b\}$ over the $1319$ distinct positive thresholds $b$, with certified weights $w_b$. So $\bar W$ is piecewise linear and $g=\bar W-0.83695P$ is convex on every linear piece.
\begin{itemize}
\item First piece: $g(0)=0$ and $g'(0)=\sum_bw_b-0.83695>4.04\times10^{-6}>0$.
\item Tail: beyond the largest threshold, $g\ge\bar W(\infty)-2(0.83695)>2.57\times10^{-5}>0$, since $P<2$.
\item Each remaining piece is split into $40$ subintervals, on which the tangent lines of $g$ at the endpoints give valid lower bounds. The smallest certified bound exceeds $2.95\times10^{-8}$, on $[0.00558303,0.00751685]$ (exact decimal endpoints).
\end{itemize}
All computations use python-flint~0.9.0 (Arb) at $256$ bits \cite{Johansson}.
\end{proof}

\begin{proposition}[computer-assisted]\label{prop:upcert}
For the explicit rational measure with $41$ layer atoms $\nu$ of the data file \texttt{inst\_do450.json}, the profile $v=\nu((\cdot,\infty))$ satisfies $V_2(0)/\int v\,dP\le0.8372914$.
\end{proposition}
\begin{proof}[Proof (certificate)]
On the grid $t_i=i/N$, $N=80000$, put $\Phi_{j}(1)=0$ and, backwards,
\[
\Phi_1(t_i)=\Phi_1(t_{i+1})+N^{-1}\Psi\bigl(\Phi_1(t_{i+1})\bigr),\qquad
\Phi_2(t_i)=\Phi_2(t_{i+1})+N^{-1}\Psi\bigl(\Phi_2(t_{i+1})-\Phi_1(t_i)\bigr),
\]
rounded up, and interpolate linearly. Both are nonincreasing. On $[t_i,t_{i+1}]$ we have $\Phi_1\ge\Phi_1(t_{i+1})$ and $\Phi_2-\Phi_1\ge\Phi_2(t_{i+1})-\Phi_1(t_i)$, so, $\Psi$ being nonincreasing, the slope conditions of Lemma~\ref{lem:general-k-verification} hold. The benchmark is $\int v\,dP=\sum_j\nu_jP(x_j)$. In ball arithmetic at $200$ bits, $\Phi_2(0)\le0.8372913936$ and $\int v\,dP\ge0.9999999999$.
\end{proof}

\section{An upper bound for the nonlinear-system value}\label{sec:theta}
The system of \cite{BPV} reads, on $[0,1)$ with $y_j(0)=1$ and $y_j(t)\to0$ as $t\uparrow1$, $\Gamma_\ell(x)=\int_x^\infty s^{\ell-1}e^{-s}\,ds$,
\[
\bigl(\Gamma_2(-\ln y_2)\bigr)' = 2 - \tfrac{1}{\theta_2} - \Gamma_3(-\ln y_2),\qquad
\bigl(\Gamma_2(-\ln y_1)\bigr)' = 2 - \Gamma_3(-\ln y_1) - \tfrac{\theta_2}{\theta_1}\bigl(2 - \Gamma_3(-\ln y_2)\bigr).
\]
Write $u_j=-\ln y_j$ and $z_j=\Gamma_2(u_j)=(1+u_j)e^{-u_j}\in(0,1]$. Also put $\Gamma_3(u)=(u^2+2u+2)e^{-u}$ and $H(z)=2-\Gamma_3(\Gamma_2^{-1}(z))$. Since $dH/dz=-\Gamma_2^{-1}(z)\le 0$, $H$ is nonincreasing, and the system reads $z_2'=H(z_2)-1/\theta_2$, $z_1'=H(z_1)-r\,f(t)$ with $r=\theta_2/\theta_1$ and $f=2-\Gamma_3(u_2)\in[0,2]$.

\begin{lemma}\label{lem:theta2}
Any solution has $\theta_2<1/2$ and $\mathcal T(\theta_2)=1$, where $\mathcal T(\theta)=\int_0^\infty ue^{-u}\,du/(1/\theta-2+\Gamma_3(u))$. The function $\mathcal T$ is increasing, and $\mathcal T(0.48300)\le 0.99990<1<1.00009\le\mathcal T(0.48302)$, so $\theta_2^\star\in(0.48300,0.48302)$.
\end{lemma}
\begin{proof}
If $\theta_2\ge 1/2$, then $z_2'\ge -\Gamma_3(u_2)$. Since $\Gamma_3(u)=(1+u)\Gamma_2(u)+e^{-u}\le(2+u)\Gamma_2(u)$, and $u_2\le 2\ln(1/z_2)$ once $u_2\ge 3$, this gives $z_2'\ge -C z_2\ln(1/z_2)$ near $0$, so $z_2$ cannot reach $0$ in finite time; the integral also diverges at $\theta=1/2$. For $\theta_2<1/2$, $z_2$ is strictly decreasing and its hitting time of $0$ is $\mathcal T(\theta_2)=1$. The enclosures use rigorous Arb quadrature on $[0,60]$ with the tail over $[60,\infty)$ bounded by $61e^{-60}/(1/\theta-2)$.
\end{proof}

\begin{proposition}[computer-assisted]\label{prop:theta}
For any solution, $r=\theta_2/\theta_1>1.37$. Consequently $\theta_1^\star+\theta_2^\star<0.48302\,(1+1/1.37)<0.835590$.
\end{proposition}
\begin{proof}
Since $H$ is nonincreasing, a continuous piecewise-linear $\phi$ with $\phi(0)\le z(0)$ and $\phi'\le H(\phi)-K(t)$ almost everywhere is a lower barrier for any solution of $z'=H(z)-\tilde K(t)$ with $\tilde K\le K$: where $\phi>z>0$, $(\phi-z)'\le H(\phi)-H(z)\le 0$. Both barriers start at $\phi_j(0)=1=z_j(0)$ and live on a grid of step $1/2000$. While $\phi_2(t_i)$ is certified positive, $\phi_2$ has slope $\min\{0,\underline H(\phi_2(t_i))-K\}$ on $[t_i,t_{i+1}]$, with $K=1/0.48300\ge 1/\theta_2^\star$ and a certified lower bound $\underline H$ evaluated only at certified points of $(0,1]$; once $\phi_2$ is not certified positive the construction stops, since $z_2\ge0$ holds trivially (the certificate records this with the sentinel value $-1$). This gives $z_2\ge\phi_2^+$ on $[0,1)$. Because $z_2$ is decreasing and $f=2-\Gamma_3(u_2)$ is a decreasing function of $z_2$, $f\le\bar f_i:=2-\Gamma_3(\overline u(\phi_2(t_{i+1})))$ on $[t_i,t_{i+1}]$ while $\phi_2(t_{i+1})$ is certified positive, where $\overline u(z)\ge\Gamma_2^{-1}(z)$ is a certified rational upper bound, and $\bar f_i=2$ afterwards. If $r\le 1.37$, the barrier $\phi_1$ built with $K=1.37\bar f_i$ satisfies $z_1\ge\phi_1\ge0.0105>0$ on $[0,1)$, contradicting $z_1(t)\to 0$. Finally $\theta_1^\star+\theta_2^\star=\theta_2^\star(1+1/r^\star)$.
\end{proof}

\begin{proof}[Proof of Theorem~\ref{thm:main}]
Existence follows from Theorem~\ref{thm:general-k-game} with $k=2$.
Proposition~\ref{prop:lowcert} and Lemma~\ref{lem:general-k-lower-transfer}
give the lower bound. Proposition~\ref{prop:upcert} and
Proposition~\ref{prop:general-k-upper-transfer} give the upper bound,
since $0.8372914<0.83730$. Proposition~\ref{prop:theta} gives
$\theta_1^\star+\theta_2^\star<0.48302(1+1/1.37)=5723787/6850000$.
Finally, exact rational subtraction gives
\[
 \frac{83695}{100000}-\frac{5723787}{6850000}
 =\frac{18641}{13700000}>0.
\]
\end{proof}
\section{Discussion and scope}\label{sec:discussion}
\paragraph{What the separation establishes.}
The cumulative formulation permits a certified construction above the nonlinear-system value. It does not determine why that particular value falls short. The loss may lie in the dual restriction used in \cite{BPV}, in the feasible construction chosen within that restriction, or in both. In particular, we do not establish suboptimality of the entire restricted pointwise programme or necessity of a nonzero dual variable $\eta$. The finite-horizon linear-programming characterisation and the lower-bound theorem of \cite{BPV} remain unaffected.

\paragraph{Computation versus proof.}
A floating-point double-oracle search used a $320$-point quantile grid, $64$ time pieces and $450$ iterations. The reported grid value $0.837050$ and profile estimate $0.837281$ are discovery diagnostics, not certified continuum bounds. Restricting the adversary to a grid can overstate a policy's off-grid guarantee. Only the explicit ball-arithmetic inequalities in Propositions~\ref{prop:lowcert}, \ref{prop:upcert} and \ref{prop:theta}, combined with the analytic transfers, justify Theorem~\ref{thm:main}. No numerical interval for $k\ge3$ is established here.

Mixtures are convenient finite representations for search and certification. Proposition~\ref{prop:general-k-occupancy-domination} shows that a bounded measurable deterministic quantile policy can dominate each finite mixture. Therefore the better certificate from the $36$-policy mixture than from the illustrative $24$-piece policy in Appendix~\ref{app:policy} is not evidence that mixing is intrinsically superior. The recovered threshold need not retain the original time mesh; attainment of the optimal value by one bounded threshold policy remains unproved.

\paragraph{Remaining questions.}
Theorem~\ref{thm:general-k-game} concerns each fixed $k$, not a budget growing with $n$. Its maximum is attained in a compact closure, not necessarily by a single implementable optimizer. It proves neither equality of the limit with the all-horizon infimum nor an exact numerical constant. Quantitative transfer bounds depend on the chosen profile or policy, and are not a uniform convergence rate over the entire continuum game. Sharper certificates for larger budgets and a characterization of extremizers are separate problems.

\paragraph{Reproducibility and provenance.}
The accompanying repository contains the rational policy mixture, layer measure, illustrative policy, certificate programs and replay documentation. Threshold-array row $j$, counted from $1$, contains the threshold for $j$ remaining slots. A layer atom $(x,w)$ specifies a jump of height $w$ in the profile; it is not an atom of an observation distribution. The upper transfer specifies the actual horizon-dependent observation law.

The analytic proof covers every fixed budget. Numerical replay covers only the supplied finite $k=2$ witnesses and the implemented arithmetic checks; it is not a formal verification of the manuscript, a proof of global numerical optimality, or an independent human peer review. The revision separates historical search outputs from certificate evidence and does not treat unsupported older numerical upper estimates as results. Automated research assistance was used in proof development, revision and internal checking. The supplied review is advisory evidence, not an external publication endorsement.

\paragraph{Declarations.}
The study uses mathematical constructions and no human-participant or personal data; ethics approval is not applicable. Data and code are prepared for the public repository \url{https://github.com/ipitchford/iid-multiunit-prophet-limits}, with the versioned archive assigned the reserved DOI above. These identifiers do not by themselves assert that public release has completed. Authorship is anonymous. No claim about undisclosed funding or competing interests is made.

\appendix
\section{An illustrative two-unit threshold policy}\label{app:policy}
\begin{proposition}[computer-assisted]\label{prop:det}
For the rational deterministic policy of Table~\ref{tab:policy} ($T=24$), $W(x)\ge0.8365\,P(x)$ for all $x\ge0$.
\end{proposition}
\begin{proof}[Certificate]
As for Proposition~\ref{prop:lowcert} with a single policy: $g'(0)>1.85\times10^{-5}>0$, the tail margin exceeds $0.0028$, and the smallest certified lower bound over the $47$ remaining pieces ($400$ subintervals each) exceeds $3.2\times10^{-5}$. These are conservative lower bounds, not equalities obtained by rounding displayed decimals.
\end{proof}

\begin{longtable}{lcc}
\caption{The certified policy ($k=2$, $T=24$). On each time piece the policy accepts an arrival with scaled tail quantile $s\le a_2$ when two slots remain, and $s\le a_1$ when one remains. All values are exact decimals.}\label{tab:policy}\\
\toprule time piece & $a_2$ (two slots) & $a_1$ (one slot)\\\midrule
$[0/24,1/24)$ & $1.170$ & $0.184$ \\
$[1/24,2/24)$ & $1.303$ & $0.123$ \\
$[2/24,3/24)$ & $1.433$ & $0.125$ \\
$[3/24,4/24)$ & $1.643$ & $0.158$ \\
$[4/24,5/24)$ & $1.781$ & $0.252$ \\
$[5/24,6/24)$ & $1.911$ & $0.306$ \\
$[6/24,7/24)$ & $2.130$ & $0.370$ \\
$[7/24,8/24)$ & $2.272$ & $0.439$ \\
$[8/24,9/24)$ & $2.413$ & $0.507$ \\
$[9/24,10/24)$ & $2.650$ & $0.581$ \\
$[10/24,11/24)$ & $2.801$ & $0.652$ \\
$[11/24,12/24)$ & $2.952$ & $0.733$ \\
$[12/24,13/24)$ & $3.224$ & $0.809$ \\
$[13/24,14/24)$ & $3.374$ & $0.889$ \\
$[14/24,15/24)$ & $3.662$ & $0.977$ \\
$[15/24,16/24)$ & $3.883$ & $1.072$ \\
$[16/24,17/24)$ & $4.086$ & $1.547$ \\
$[17/24,18/24)$ & $4.483$ & $2.040$ \\
$[18/24,19/24)$ & $4.723$ & $2.540$ \\
$[19/24,20/24)$ & $5.297$ & $3.068$ \\
$[20/24,21/24)$ & $5.620$ & $3.593$ \\
$[21/24,22/24)$ & $6.009$ & $4.234$ \\
$[22/24,23/24)$ & $8.273$ & $5.010$ \\
$[23/24,24/24)$ & $10.602$ & $6.587$ \\
\bottomrule
\end{longtable}

\begin{thebibliography}{99}
\bibitem{BPV} J.~Brustle, S.~Perez-Salazar, V.~Verdugo. Splitting guarantees for prophet inequalities via nonlinear systems. \emph{Math.\ Oper.\ Res.} 51(2) (2026) 877--904 (online 17 April 2025), doi:10.1287/moor.2024.0413; arXiv:2406.17767v2.
\bibitem{Brustle26} J.~Br\"ustle. \emph{Optimisation for prophet inequalities}. PhD thesis, London School of Economics (2025), doi:10.21953/lse.00004808.
\bibitem{CPV26} E.~Cruz-Ossa, S.~Perez-Salazar, V.~Verdugo. Competition versus complexity in multiple-selection prophet inequalities. arXiv:2602.20398 (2026).
\bibitem{CFHOV21} J.~Correa, P.~Foncea, R.~Hoeksma, T.~Oosterwijk, T.~Vredeveld. Posted price mechanisms and optimal threshold strategies for random arrivals. \emph{Math.\ Oper.\ Res.} 46 (2021) 1452--1478.
\bibitem{Fan} K.~Fan. Minimax theorems. \emph{Proc.\ Natl.\ Acad.\ Sci.\ USA} 39 (1953) 42--47.
\bibitem{GMZ} A.~Goldenshluger, Y.~Malinovsky, A.~Zeevi. Optimal single-threshold stopping rules and sharp prophet inequalities. \emph{J.\ Appl.\ Probab.} (2026), First View, 1--28 (online 27 July 2026), doi:10.1017/jpr.2026.10117; arXiv:2404.12949.
\bibitem{HK82} T.~P.~Hill, R.~P.~Kertz. Comparisons of stop rule and supremum expectations of i.i.d.\ random variables. \emph{Ann.\ Probab.} 10 (1982) 336--345.
\bibitem{Kertz86} R.~P.~Kertz. Stop rule and supremum expectations of i.i.d.\ random variables: a complete comparison by conjugate duality. \emph{J.\ Multivariate Anal.} 19 (1986) 88--112.
\bibitem{JMZ} J.~Jiang, W.~Ma, J.~Zhang. Tightness without counterexamples: a new approach and new results for prophet inequalities. \emph{Proc.\ EC} 2023; \emph{Math.\ Oper.\ Res.}\ 51(2) (2026) 956--987 (online 2025), doi:10.1287/moor.2023.0221; arXiv:2205.00588v5.
\bibitem{Johansson} F.~Johansson. Arb: efficient arbitrary-precision midpoint-radius interval arithmetic. \emph{IEEE Trans.\ Comput.} 66 (2017) 1281--1292.
\bibitem{KM26} J.~Kong, K.~Murthy. Multiunit i.i.d.\ prophet inequalities via extreme value asymptotics. arXiv:2602.18756v2 (2026).
\bibitem{PST} S.~Perez-Salazar, M.~Singh, A.~Toriello. The IID prophet inequality with limited flexibility. \emph{Math.\ Oper.\ Res.} 51(1) (2026) 218--254 (online 2025), doi:10.1287/moor.2022.0345; arXiv:2210.05634.
\bibitem{PSV} S.~Perez-Salazar, V.~Verdugo. Optimal guarantees for online selection over time. arXiv:2408.11224 (2024).
\bibitem{LeCam} L.~Le~Cam. An approximation theorem for the Poisson binomial distribution. \emph{Pacific J.\ Math.} 10 (1960) 1181--1197.
\bibitem{LMPPV} P.~Loiseau, M.~Molina, V.~Perchet, S.~Perez-Salazar, V.~Verdugo. Kernel methods for refined prophet inequalities. arXiv:2608.08662.
\bibitem{MGLP} M.~Molina, N.~Gast, P.~Loiseau, V.~Perchet. Prophet inequalities: competing with the top $\ell$ items is easy. \emph{Proc.\ SODA} (2025) 1270--1307, doi:10.1137/1.9781611978322.38; arXiv:2408.07616v2.
\end{thebibliography}
\end{document}
